\section{Related Work}
\label{sec:related}

\paragraph{Compositional reasoning frameworks.}
The classic foundations are Hoare's axiomatic semantics~\cite{hoare1969axiomatic}
and Dijkstra's weakest-precondition calculus~\cite{dijkstra_guarded_1975,dijkstra1976discipline}.
Cousot \& Cousot's abstract interpretation~\cite{cousot1977abstract} provides
the general framework within which our second-pass refinement sits: the
isolated pass is a sound but conservative approximation of the
fixpoint; our refinement narrows it via callee summaries. Boogie's
verification-condition generator~\cite{barnett2005weakest} performs the
full WP transformation that our scaffold approximates; Boogie targets
verification rather than the cheap-and-rich JML emission this work
serves.

\paragraph{Procedure summaries and interprocedural analysis.}
Calcagno et al.'s Infer~\cite{calcagno2015infer} computes compositional
summaries via bi-abduction in separation logic; the architecture is
genuinely compositional from the outset and does not need a
``second pass'' because the first pass is already
compositional. Houdini~\cite{flanagan2001houdini} conjectures
candidate annotations and prunes them via ESC/Java calls; the
loop is closer to test-and-refute than to constructive
compositional inference but the spirit is similar. Daikon's likely
invariants~\cite{ernst2007daikon} are observed dynamically per
program point and do not have an analogue of a second pass; the
question we address (does single-pass propagation capture
compositional information) is specific to static syntactic
inferrers.

\paragraph{Specification inference for JML.}
The immediate predecessor is~\cite{edmonds2026inference}, which
introduced the heuristic JML inferrer we extend, and
~\cite{edmonds2026embedding}, which addressed the distribution
problem (embedding inferred specs in compiled bytecode and OpenAPI
descriptions). PreInfer~\cite{preinfer2017} uses dynamic symbolic
execution to infer preconditions; its propagation strategy is
trace-driven rather than syntactic and is not directly comparable.
Jdoctor~\cite{jdoctor2018} extracts exception preconditions from
Javadoc; its scope is narrower than ours but its precondition
synthesis is purely textual and does not propagate.

\paragraph{Polymorphic dispatch in static analysis.}
Tip et al.~\cite{tip2002practical} address polymorphic dispatch via
class-hierarchy analysis (CHA) and rapid-type analysis (RTA), the
techniques the inferrer's \texttt{CallGraph} uses to enumerate
candidate callees. Our disjunction-over-candidates is the static
equivalent of summarising the runtime-may-dispatch set. The
trade-off is the standard CHA precision/conservatism trade-off:
disjunction over a large set of candidates produces preconditions
that are easy to satisfy and therefore less informative.

\paragraph{Mutation testing as a downstream evaluation.}
PIT~\cite{coles2016pitest} is the mutation-testing pipeline the
predecessor work uses to evaluate downstream spec quality. The
present article does not run a mutation experiment; we report only
the structural-extension claim. The follow-up evaluation will run
PIT against test suites generated from the isolated baseline and
from the compositionally-refined specs to ask whether the additional
clauses translate into more bugs caught at the suite level.
